\begin{lemma}[Restricted augmentation of a good hierarchy]
Let $G$ be a finite loopless multigraph retaining individual edge
identities, let $e_0$ be distinguished, and let $[k]$ be the color set.
Let $\mathcal E$ be a positive-rung ETT with generating coloring
$\varphi$, maximum property, terminal ladder tree $T_n$, and final
extension kind $\Theta_n$.  Let $\sigma$ be a $\varphi\bmod T_n$
coloring.  Suppose that an ordered tree $J$ carries a corresponding ETT
$\mathcal E_J$ for $(\sigma,T_n)$, a hierarchy base $B_J$, a hierarchy
$K_J$, and a good-hierarchy witness $\mathcal G_J$.  Put
\[
 d=K_J.\mathit{dividerCount},\qquad
 A=K_J.\mathit{starred}(d),\qquad D=D^{K_J}_{n,d},
\]
and write $\Gamma_\eta=\mathcal G_J.\Gamma^d_\eta$.

For an exact ordered extension $W$ of $A$, let $C_W(\eta)$ be the set of
colors on the exact displayed edges of $W$ whose positions are at least
$|A|$ and strictly below both $|W|$ and the totalized first position in
$W$ missing $\eta$.  Then there are a natural number $m$, exact ordered
trees
\[
                  U_0=J,U_1,\ldots,U_m=U,
\]
and a terminal tree $U$ such that all of the following hold.
\begin{enumerate}
\item For every $r<m$, $U_{r+1}$ is one
  $\sigma$-\noderef{TashkinovAugmentation} of $U_r$, and
  $\Gamma_\eta\cap C_{U_{r+1}}(\eta)=\varnothing$ for every $\eta\in D$.
  Moreover $|U_r|=|J|+r$, every $U_r$ is an exact ordered initial segment
  of every later $U_s$, and $V(U_r)$ is a proper subset of
  $V(U_{r+1})$.  In particular, the displayed chain is an actual
  $\sigma$-\noderef{TashkinovAugmentationSequence} from $J$ to $U$.
\item The terminal tree still satisfies
  $\Gamma_\eta\cap C_U(\eta)=\varnothing$ for every $\eta\in D$, and it
  is maximal for this property: no one-step Tashkinov augmentation $W$ of
  $U$ satisfies all these same fixed last-index avoidance conditions.
\item The tree $U$ carries a corresponding ETT $\mathcal E_U$ for
  $(\sigma,T_n)$ with maximum property, a reconstructed hierarchy base
  $B_U$, and a hierarchy $K_U$ with $d$ dividers.  The core, closure, and
  numerical supporting-vertex position of $B_U$ equal those of $B_J$.
  After the canonical identification of divider indices, all nonterminal
  stages, starred stages, indexed stored sets, and stage-missing sets of
  $K_U$ equal those of $K_J$, and the numerical divider positions agree.
  There is a good witness $\mathcal G_U$ whose every indexed Gamma pair
  equals the corresponding pair of $\mathcal G_J$.
\item The terminal stored set satisfies
  \[
                  D^{K_U}_{n,d+1}\subseteq D^{K_J}_{n,d},
  \]
  and $U$ is
  $\bigl(\bigcup_{\eta\in D^{K_U}_{n,d+1}}
  \mathcal G_U.\Gamma^d_\eta\bigr)^-$-closed under $\sigma$.
\end{enumerate}
Here every initial-segment assertion and every equality of stages concerns
the complete ordered edge-valued tree-sequences, not merely their vertex
sets.
\end{lemma}
\begin{proof}
Write $D=D^{K_J}_{n,d}$ and $\Gamma_\eta=\Gamma^d_\eta$ throughout.
By \noderef{HierarchySegmentColors}, the set denoted $C_W(\eta)$ in the
statement is exactly the last-stage segment-color expression obtained by
using the fixed starred tree $A$, the next stage $W$, and
\noderef{HierarchyFirstMissingPosition}$(\sigma,W,\eta)$.  In particular,
the last avoidance clause of \noderef{GoodHierarchy} for $K_J$ says
\[
                    \Gamma_\eta\cap C_J(\eta)=\varnothing
                    \quad(\eta\in D).                 \tag{1}
\]

We first make the one-edge construction used by the restricted loop
completely explicit.  Let $X$ be a current tree and suppose that an exact
edge identity $f$ belongs to
$\noderef{BoundaryColorEdges}(G,\sigma,V(X),c)$ for a color
$c\in\noderef{TreeMissingColors}(\sigma,X)$.  Unfolding
\noderef{BoundaryColorEdges} and \noderef{BoundaryEdges} supplies vertices
$x\in V(X)$ and $y\notin V(X)$ with
$G.\mathit{ends}(f)=\{x,y\}$.  Define $X+f$ to have length $|X|+1$; copy
all old vertex and edge entries, put $y$ in the new last vertex position,
and put the individual identity $f$ in the new last edge position.

We verify every field of this new \noderef{TreeSequence}.  Positivity
follows from the positivity of $X$.  For vertex injectivity, two old
indices reduce to injectivity in $X$, while equality involving the new
last entry would put $y$ in $V(X)$, contrary to its choice.  For edge
injectivity, two old indices again reduce to injectivity in $X$.  The new
identity $f$ cannot equal an old displayed edge: the attachment field of
$X$ puts both endpoints of every old edge in $V(X)$, whereas $f$ has the
endpoint $y\notin V(X)$.  The first-edge equation is copied because the
old positive tree already contains its distinguished first edge.  Every
old attachment equation is copied.  For the new last edge, choose the old
position occupied by $x$; symmetry of the unordered endpoint pair rewrites
$\{x,y\}$ as $\{y,x\}$, giving precisely the required attachment equation.
Thus the constructed object is an actual tree-sequence with fresh vertex
and edge entries, not merely a vertex-set extension.

Its length equation, entrywise old vertex equalities, and entrywise old
edge equalities give
$\noderef{TreeSequenceInitialSegment}(X,X+f)$; the same data and the exact
last-edge equality give \noderef{TreeSequenceAdjoinEdge}$(X,f,X+f)$.
Together with $\sigma(f)=c$ and $c\in\overline\sigma(X)$, these are all
fields of \noderef{TashkinovAugmentation}$(\sigma,X,X+f)$.  Directly from
the copied entries, the fresh last vertex, and
\noderef{TreeSequenceVertexSet},
\[
                V(X+f)=V(X)\cup\{y\},\qquad V(X)\subsetneq V(X+f).
                                                               \tag{2}
\]

Call a finite chain admissible if it starts at $J$, every successor is a
one-step $\sigma$-Tashkinov augmentation, and at every successor $W$ it
satisfies
\[
                    \Gamma_\eta\cap C_W(\eta)=\varnothing
                    \quad(\eta\in D).                 \tag{3}
\]
The length-zero chain is admissible by (1).  At every successor the
definition of \noderef{TashkinovAugmentation} gives one-more length and an
exact prefix.  Thus there is exactly one new vertex position.  Injectivity
of the successor's vertex array makes its entry different from every old
entry, while exact prefix equality identifies all the old entries; this
directly gives (2) at every step.  The identical argument for its edge
array identifies one fresh edge identity as well.  Since all these
vertex sets are subsets of the fixed finite type $V$, an admissible chain
has at most $|V|-|V(J)|$ successors.  The nonempty finite set of attainable
chain lengths therefore has a largest member.  Choose an admissible chain
$U_0,\ldots,U_m$ of largest length and put $U=U_m$.  If an admissible
one-step successor of $U$ existed, appending it would give an admissible
chain of length $m+1$, a contradiction.  This proves the asserted
maximality.  The asserted avoidance at the endpoint also holds: it is (1)
when $m=0$, and it is the defining successor condition (3) at $r=m-1$
when $m>0$.

Induction on $s-r$, using the exact initial segment at each successor and
entrywise composition, proves that $U_r$ is an exact
\noderef{TreeSequenceInitialSegment} of $U_s$ whenever $r\leq s$.
The one-more length equation similarly gives $|U_r|=|J|+r$.  The same
induction, now on the relation defined in
\noderef{TashkinovAugmentationSequence}, proves that the whole displayed
chain is a Tashkinov augmentation sequence from $J$ to $U$.  Thus every
chain and exact-prefix assertion in items 1 and 2 has been obtained.

We next construct a corresponding ETT on every $U_r$; retaining all of
them during the construction will also justify the terminal certificate.
Let $n$ be the last rung, let $T_n$ be the common terminal ladder tree,
and let $Z$ be the final closure tree in the replacement series displayed
by the given \noderef{CorrespondingExtendedTashkinovTree} certificate on
$J$.  That certificate and \noderef{ExtendedTashkinovTree} give the exact
segments
\[
                         T_n\preccurlyeq J\preccurlyeq Z,
                         \qquad |T_n|<|J|.              \tag{4}
\]
Because the rung is positive, the certified step at its last rung is one
of the SE, PE, or RE constructors of \noderef{TashkinovSeriesStep}, not
the initial constructor: the certification function itself is the field of
\noderef{TashkinovSeries} carried by the replacement ETT.

In the PE case, that constructor contains a
\noderef{TashkinovClosure} from $T_n$ to $Z$.  Its first field is a
$\sigma$-Tashkinov augmentation sequence.  Apply
\noderef{TashkinovAugmentationSequenceRestrict} to this sequence and the
two exact segments in (4); it gives an augmentation sequence from $T_n$
to $J$.  Appending the accumulated chain from $J$ to $U_r$ gives one from
$T_n$ to $U_r$.

In the SE and RE cases, the terminal constructor displays an individual
edge $f_n$, an exact one-edge seed $S=T_n+f_n$, and a
\noderef{TashkinovClosure} from $S$ to $Z$.  The seed has length
$|T_n|+1\leq|J|$ by (4).  Both $S$ and $J$ are exact prefixes of $Z$: for
$S$ this follows by induction along the closure's augmentation sequence,
and for $J$ it is (4).  Comparing their entries up to $|S|$ therefore gives
$S\preccurlyeq J$.  Apply
\noderef{TashkinovAugmentationSequenceRestrict} to the closure route
$S\leadsto Z$ and the segments $S\preccurlyeq J\preccurlyeq Z$; this
gives $S\leadsto J$, which again appends to $J\leadsto U_r$.
The SE or RE constructor also says that its replacement connecting-edge
set is exactly $\{f_n\}$.  Connecting-edge agreement in
\noderef{CorrespondingExtendedTashkinovTree} therefore identifies the
base series' exact set with the same singleton $\{f_n\}$.  We have now
checked the appropriate PE or SE/RE antecedent, including the exact
connecting-edge identity, in the universal clause of
\noderef{ModColoring}.

Equation (4) and $J\preccurlyeq U_r$ give
$T_n\preccurlyeq U_r$ and $|T_n|<|U_r|$.  The universal mod-coloring clause
therefore supplies a corresponding ETT $\mathcal E_{U_r}$ for every
$U_r$.  In particular it supplies $\mathcal E_U$.  Apply
\noderef{MaximumPropertyTransport} to the given
\noderef{MaximumProperty} of $\mathcal E$: the resulting ETT
$\mathcal E_U$ has maximum property.

We now reconstruct the hierarchy base $B_U$, checking its fields rather
than transporting the old structure across unlike dependent types.  The
two corresponding certificates identify their terminal ladder trees,
raw stored sets, exact connecting-edge sets, and extension kinds with the
same data of $\mathcal E$.  Transport the numerical support position of
$B_J$ through the equality of the terminal ladder trees, and retain its
positivity.  Retain also the old closure tree and core.  The terminal
missing-color compatibility field follows from the identical ladder tree
and the common coloring $\sigma$.

If $\Theta_n$ is not PE, the nonparallel field of
\noderef{HierarchyBase} identifies the old core with the old terminal
ladder tree.  Ladder-tree agreement identifies this with the new terminal
ladder tree, proving the new nonparallel field.  The parallel implication
has a false antecedent in this case.

Suppose instead that $\Theta_n$ is PE.  Unpack the terminal PE constructors
of the certified replacement series for $\mathcal E_J$ and
$\mathcal E_U$.  Each constructor contains a
\noderef{MaximumDefectiveChoice}.  Connecting-edge agreement and the two
singleton constructor equations give equality of the individual chosen
edge identities $f_J=f_U$.  For each choice,
\noderef{DefectiveVertex} and \noderef{MaximumDefectiveChoice} put its
chosen vertex inside the common $V(T_n)$, put its other endpoint outside
$V(T_n)$, give \noderef{Incident} of the chosen vertex with that edge, and
give the exact unordered endpoint equation.  Equality of the edge
identities makes the two unordered pairs equal.  The inside endpoint of
one pair cannot equal the outside endpoint of the other, so the
inside/inside alternatives are forced; hence the two support vertices are
equal.  The old parallel field now identifies that common support vertex
with the vertex at the retained numerical support position.

The closure stored in $B_J$ is already a
\noderef{TashkinovClosure} under $\sigma$ of the prefix of $T_n$ at that
position.  Equality of $T_n$, of the support position, and the fieldwise
definition of \noderef{TreeSequencePrefix} identify the new prefix with
the old one, so the same closure proof applies.  Finally the old
\noderef{OrderedTreeUnion} witness depends only on this common $T_n$, the
retained closure, and the retained core; reuse it verbatim.  Support
positivity is unchanged, and the terminal missing-color equality was
checked above.  These arguments construct every field of $B_U$ and prove
the three advertised data equalities.

Construct $K_U$ with divider count $d$, with stages
\[
 K_U(0)=K_J(0),\ldots,K_U(d)=K_J(d),\qquad K_U(d+1)=U.
                                                               \tag{5}
\]
The base-stage equation copies.  All nesting and strictness fields before
the last one copy.  The last follows by composing the old last segment
$K_J(d)\preccurlyeq J$ with $J\preccurlyeq U$; its strictness already
holds before that composition.  Embed each old divider position
$p<|J|+1$ into $\operatorname{Fin}(|U|+1)$ using $|J|\leq|U|$.  Numerical
values, positivity, and strict increase are consequently unchanged.

For a divider position $p$, the old divider-stage field identifies the
relevant stage with \noderef{TreeSequencePrefix}$(J,p)$.  Since
$J\preccurlyeq U$, the prefixes of $J$ and $U$ through $p$ have the same
length and the same vertex and edge entries.  Their attachment and
injectivity fields are propositions, so proof irrelevance identifies the
two tree-sequences.  This proves every new divider-stage equation
entrywise.  If $d>0$, the first ordinary stage in (5) is unchanged, and
the old core-inside-first field copies.  If $d=0$, the first ordinary stage
is now the terminal $U$; compose the old relation
$B_J.\mathit{core}\preccurlyeq J$ with $J\preccurlyeq U$, and use equality
of the reconstructed cores.  This separately proves core-inside-first in
the zero-divider case.  All fields of \noderef{TashkinovHierarchy} have
now been checked.  Equation (5), the definition of
\noderef{HierarchyStarredStage}, and the direct definitions of
\noderef{HierarchyStoredColorsAt} and
\noderef{HierarchyStageMissingColors} give all indexed stage, starred,
stored-set, and stage-missing equalities asserted in item 3.

Define $\mathcal G_U.\Gamma^j_\eta=\mathcal G_J.\Gamma^j_\eta$ under
these index identifications.  We check every field of
\noderef{GoodHierarchy}.  Pair cardinality and pairwise disjointness copy
using indexed stored-set equality.  Every unique-refresh clause has index
$i<d$; its two ordinary stages, preceding starred stage, vertex-set
difference, stored sets, and all Gamma pairs are unchanged, so its old
unique witness and uniqueness clause copy.  The PE base clause uses only
the common terminal series data, reconstructed equal core and closure,
zeroth stored set, zeroth Gamma pairs, and their $\sigma$-missing sets, so
both its non-minus-closed and cardinality assertions copy.  Every internal
minus-closed clause likewise concerns an unchanged nonterminal stage,
stored set, and preceding Gamma union, and therefore copies.

For avoidance at an index $j<d$, the starred stage and next ordinary stage
are unchanged, so \noderef{HierarchySegmentColors} is unchanged and the
old avoidance clause copies.  At $j=d$, its stage-missing set is unchanged,
and the terminal admissibility condition (3) at $U$ says exactly that the
old Gamma pair is disjoint from the new last segment-color set.  To check
that this includes all totalized cases, unfold
\noderef{HierarchyFirstMissingPosition}.  If $U$ has an $\eta$-missing
vertex, the least such position is used.  If it has none, the fallback is
the last vertex position $|U|$, so the strict cutoff in $C_U(\eta)$ still
includes precisely every eligible displayed edge and no nonexistent one.
Thus (3) is exactly the required avoidance clause in both the genuine and
no-missing branches.  This completes the construction of $\mathcal G_U$
and proves all indexed Gamma equalities.

It remains to prove the two terminal assertions.  We do the stored-set
comparison directly.  Unfold \noderef{TashkinovStoredColors},
\noderef{HierarchyStoredColorsAt}, and
\noderef{HierarchyTerminalStoredColors}.  The corresponding certificates
identify every raw set $S_h$ and hence their union $R$; the terminal ladder
trees and their $\sigma$-missing sets agree as well.  Thus the terminal
Tashkinov stored set on either side is
$R\setminus\overline\sigma(T_n)$.  Exact prefix composition gives
$T_n\preccurlyeq A\preccurlyeq J\preccurlyeq U$, and hence
$\overline\sigma(T_n)\subseteq\overline\sigma(A)
\subseteq\overline\sigma(U)$ by the union definition in
\noderef{TreeMissingColors}.  Consequently the iterated differences in
the two hierarchy definitions reduce, element by element, as follows:
\[
 (R\setminus\overline\sigma(T_n))
       \setminus\overline\sigma(U)=R\setminus\overline\sigma(U),
 \qquad
 (R\setminus\overline\sigma(T_n))
       \setminus\overline\sigma(A)=R\setminus\overline\sigma(A).
\]
Therefore
\[
 D^{K_U}_{n,d+1}=R\setminus\overline\sigma(U),\qquad
 D^{K_J}_{n,d}=R\setminus\overline\sigma(A).             \tag{6}
\]
Taking differences from the same $R$ in (6) proves
$D^{K_U}_{n,d+1}\subseteq D^{K_J}_{n,d}$.  Notice that this is a direct
raw-set and missing-set calculation; no adjacent-stage antitonicity result
is being used.

Finally suppose, contrary to the desired
\noderef{TreeColorSetMinusClosed} certificate, that a color
\[
 c\in\overline\sigma(U)\setminus
       \bigcup_{\eta\in D^{K_U}_{n,d+1}}\Gamma_\eta       \tag{7}
\]
has a $c$-colored boundary edge identity $f$.  Apply the fieldwise
one-edge construction above to obtain the exact Tashkinov augmentation
$W=U+f$.  We show that $W$ obeys (3), contradicting maximality.

Fix $\eta\in D$.  If $\eta\in D^{K_U}_{n,d+1}$, (7) gives
$c\notin\Gamma_\eta$.  Every old edge in $C_W(\eta)$ already lies in
$C_U(\eta)$, while the only possible new member is the color $c$ of the
new last edge; hence the old admissibility and $c\notin\Gamma_\eta$ prove
disjointness.  This reasoning covers both possibilities for the totalized
first-missing position.  If the new vertex is the first vertex missing
$\eta$, its position is the new last vertex and the new edge lies below
that position.  If no vertex of $W$ misses $\eta$, the fallback is the
same new last vertex position and the conclusion is identical.

If instead $\eta\notin D^{K_U}_{n,d+1}$, equations (6) and $\eta\in D$
show that $\eta$ is missing on $U$.  Its genuine first missing position is
therefore at most $|U|$ and is unchanged in the exact extension $W$.
The new edge has position $|U|$, so it is not strictly below that first
missing position.  Consequently $C_W(\eta)=C_U(\eta)$, and old
admissibility again gives disjointness.  Thus $W$ satisfies every fixed
last-Gamma avoidance condition, contradicting the maximality of $U$.

No color satisfying (7) can therefore have a colored boundary edge.
Unfolding \noderef{TreeColorSetClosed} and
\noderef{TreeColorSetMinusClosed} says exactly that $U$ is
$\bigl(\bigcup_{\eta\in D^{K_U}_{n,d+1}}\Gamma_\eta\bigr)^-$-closed.
Together with the reconstructed ETT, base, hierarchy, and good witness,
this proves every assertion of the lemma.
\end{proof}
